\section{Research Questions and Portability Measures}\label{sec:study-design}\label{sec:rq}
We preserve the original six questions, including those not answered by
adequate evidence in this study:
\begin{description}
  \item[RQ1.] How stable are scheduler rankings across independent
    workload sources?
  \item[RQ2.] How stable are scheduler rankings under temporal, provider,
    and domain shifts?
  \item[RQ3.] To what extent do rankings obtained on synthetic stress
    workloads transfer to rankings on independent real-trace-derived
    workloads?
  \item[RQ4.] Which source-native observable workload characteristics are
    associated with cross-dis\-tribution scheduler rank reversals?
  \item[RQ5.] How many independent workload windows are required before a
    comparative scheduler ranking becomes statistically stable?
  \item[RQ6.] Do representative simulated scheduler rankings and
    cross-workload rank reversals reproduce on a real serving engine?
\end{description}
Here ``independent sources'' means separately released artifacts, not proven
independence of providers or of every request. RQ1/RQ2 are the primary
cross-source and temporal questions. RQ3 remains an engineering pilot,
RQ4 a descriptive association analysis, RQ5 finite-reference recovery,
and RQ6 unresolved as a faithful native-system comparison. Operating-region
sensitivity is secondary. Cross-metric analysis is explicitly post-campaign;
its prominence in this manuscript does not make it retrospectively preregistered.

\subsection{Comparison object and statistical unit}\label{sec:unit}
Let $R(W,M,L)$ denote the ordering of a fixed policy panel under workload
source or window set $W$, metric $M$, and operating region $L$. Policy
outcomes are paired within the same window and region. For source-level
rankings, a policy's defined outcomes are averaged over windows, after the
two identical deterministic execution passes are collapsed. Comparisons
between sources use different window populations; they are not paired
request-by-request across sources. The window is the resampling unit.
Non-overlap and deterministic sampling do not prove statistical independence
of windows drawn from a shared production stream.

Tau-b is computed from the policy-level aggregate values, preserving equal
values as ties. Top-$k$ summaries instead sort by value and then by policy
identifier. In particular, the reported top-1 agreement is agreement of
this deterministic tie-breaking rule, not overlap of all policies tied
for the optimum. Exact-order recovery uses the same distinction. These
conventions are retained for reproducibility and bound the interpretation
of preferred-policy stability.
\subsection{Metrics}
\label{sec:metrics}

Three metrics are always defined for every completed cell: completion
fraction, weighted completion fraction, and Arrival-Normalized Weighted Goodput (ANWG, the primary metric). We formally define ANWG here because it is this benchmark's headline measure of system-level utility. For a set of arriving requests $R$, let $w_i \ge 1$ denote the priority weight assigned to request $i \in R$. Let $C \subseteq R$ be the subset of requests that successfully complete execution, and let $S \subseteq C$ be the subset of completed requests that also meet their latency deadline (SLO), with completion at or before the deadline. ANWG is defined as the priority-weighted sum of successfully attained deadlines divided by the total priority weight of all arrivals:
\begin{equation}
\text{ANWG} = \frac{\sum_{i \in S} w_i}{\sum_{i \in R} w_i}
\end{equation}
Because the denominator is the sum over all \emph{arrivals} rather than completions, rejected, dropped, unfinished, and SLO-missed requests all receive zero numerator credit while remaining in the denominator penalty. ANWG is always defined (the denominator is zero only if a workload window has no arrivals, which is prevented by construction); therefore, undefined-metric exclusions for ANWG are always exactly zero. A second group of metrics is defined
only conditional on at least one request completing: SLO-violation rate,
weighted goodput, mean/p95/p99 latency, and request/token throughput are
each \texttt{NaN} if and only if completion fraction is exactly zero for
that cell. Mean/p95 time-to-first-token (TTFT) is subject to a stricter,
separately checked precondition (at least one completed request recorded a
first-token time), because TTFT can be undefined even when completion
fraction is positive. No undefined conditional-metric value is ever imputed
as 0 or 1, in either direction, and never after inspecting which policy or
source it would affect. Where a policy's value is undefined for a given
(window, region, metric) condition, that policy is excluded from the
ranking-derived statistic (Kendall's tau-b, Spearman's rho, top-$k$
overlap, pairwise reversal detection) for that specific condition and
metric only -- never scored as a loss or a tie, and never causing the whole
condition to be dropped for policies that did complete. Every
ranking-derived table reports its exclusion count
(\texttt{n\_conditions\_excluded\_for\_undefined\_metric}) alongside the
statistic itself, since a source-concentrated exclusion rate is itself a
reportable finding under RQ4, not a footnote.

\subsection{Ranking-Uncertainty Quantification and Reversal Classification}
\label{sec:uncertainty}

Because the statistical unit of analysis is the workload window
(\S\ref{sec:unit}), ranking-uncertainty is quantified by a whole-window
block bootstrap: windows, not individual requests, are resampled with
replacement (preserving each window's internal request sequence intact),
2{,}000 resamples are drawn per condition, and every reported confidence
interval is a 95\% percentile interval over those resamples, using a
deterministic, recorded seed so every interval is exactly reproducible from
the frozen outcome matrix alone. This block structure exists because the
requests within a window are not independent observations of scheduler
behavior -- they share one arrival process, one prompt/output-length
mixture, and one operating region -- so resampling at the request level
would understate uncertainty; the window is the design's resampling unit. Residual temporal dependence between windows may still affect interval coverage; it is not removed merely by treating requests within a window as a block.

Two further checks calibrate how the pairwise reversal statistics reported
in \S\ref{sec:reversals} are interpreted, and it is important not to
conflate them. The frozen output contains 66 Friedman omnibus records (11 metrics by six regions). Earlier manuscript text incorrectly described the omnibus test as unimplemented; this draft corrects that statement. The principal reversal findings below use the separately recorded pairwise procedure.
Separately, because the panel and the (source, region, metric) grid
together generate many pairwise comparisons, every pairwise reversal
test's two-sided bootstrap sign-test $p$-value -- the intersection-union
combination $\max(p_x, p_y)$ of the two compared conditions' sign-test
$p$-values, read off the same resamples that back the ordinary CI rule
below -- is additionally checked against a Benjamini--Hochberg
false-discovery-rate (FDR) correction at $q = 0.05$, applied per family, where a
family is all pairwise policy comparisons within a single (metric,
load-region) combination. This produces a separate, multiplicity-corrected
\texttt{supported\_after\_fdr} flag on every record. In the frozen
benchmark artifact, all $36$ records that the five-class taxonomy below
independently classifies as \texttt{SUPPORTED\_PRACTICAL\_REVERSAL} also
satisfy \texttt{supported\_after\_fdr} -- the FDR check removes no headline
reversal in this dataset -- but the two criteria are computed
independently from each other and are not guaranteed to coincide in
general.

A numerical sign change in a point estimate is not, by itself, a meaningful
scheduler reversal from a systems standpoint: two policies can swap
point-estimate order by a magnitude far smaller than the benchmark's own
chosen practical margin, in which case treating it as a reversal would overstate
what the benchmark actually supports. The 10\% winning-margin threshold
used throughout this test was fixed prospectively, before any
reversal was observed, by reusing this project's own earlier
discriminability-pilot criterion (\texttt{\_relative\_range\_qualifies},
\S\ref{sec:stage0}) rather than adopting a threshold chosen after seeing
these results or asserted as a domain-standard cutoff; it is an
operational convention for separating materially sized ordering changes
from microscopic sign changes, not a claim that 10\% is a universal
practical-significance threshold for LLM serving, and every count labeled
``practical'' in this manuscript is conditional on it. Every pairwise (policy A, policy B)
comparison at a given (source, region, metric) condition is therefore
classified into exactly one of five preregistered classes, combining the
practical-margin test with the \emph{ordinary} (not FDR-adjusted) 95\%
percentile bootstrap-CI test defined above; the
FDR check described above is a separate, additional multiplicity
safeguard applied on top of this classification for headline reporting,
not an input the classification itself depends on:

\begin{description}
  \item[\texttt{UNDEFINED\_UNESTIMABLE}] At least one of the two policies'
    metric values is undefined for this condition (\S\ref{sec:metrics});
    no direction is estimated or claimed.
  \item[\texttt{STABLE\_NO\_SIGN\_CHANGE}] The two policies' relative order
    agrees across the compared conditions; no reversal occurred.
  \item[\texttt{MICROSCOPIC\_SIGN\_CHANGE}] The point-estimate order
    reverses, but the winning margin fails to exceed 10\% of the losing
    policy's value in \emph{at least one} of the two compared conditions --
    this includes both the case where the margin is negligible on both
    sides and the asymmetric case where it exceeds 10\% on one side but not
    the other. Either way the sign flip is reported here as numerically
    real but practically negligible or unreliable, never pooled with the
    classes below.
  \item[\texttt{UNSUPPORTED\_SIGN\_CHANGE\_WIDE\_CI}] The point-estimate
    order reverses and the winning margin exceeds 10\% of the losing
    policy's value in \emph{both} compared conditions, but the
    \emph{ordinary} (not FDR-adjusted) 95\% bootstrap interval for the
    difference fails to
    exclude zero in \emph{at least one} of the two conditions -- the
    apparent reversal is not statistically distinguishable from bootstrap
    noise, at the preregistered confidence level, in at least one of the
    two conditions being compared.
  \item[\texttt{SUPPORTED\_PRACTICAL\_REVERSAL}] The point-estimate order
    reverses \emph{and} the winning margin exceeds 10\% of the losing
    policy's value in both compared conditions \emph{and} the
    \emph{ordinary} (not FDR-adjusted) 95\% bootstrap interval for the
    direction excludes zero
    in both conditions. Only this class is reported as a headline reversal
    in \S\ref{sec:reversals}; the other four classes are always reported
    alongside it, never omitted as a residual category. Every record in
    this class additionally satisfies the separate FDR-corrected
    \texttt{supported\_after\_fdr} check described above in the frozen
    benchmark artifact.
\end{description}

These five classes are mutually exclusive and exhaustive by construction:
\texttt{classify\_pairwise\_reversal} in the frozen analysis code evaluates
undefined-value, sign-change, margin, and confidence-interval conditions as
a strict if/elif/else chain with an unconditional final branch, so every
possible (policy A, policy B, condition) input reaches exactly one class,
including the asymmetric-margin and single-sided-wide-CI cases described
above; no already-computed reversal record is at risk of being silently
dropped or misclassified by a gap in this taxonomy.


\subsection{Aggregation, eligibility and reproducible interpretation}
Write $y_{w,p,m}$ for the defined metric value of policy $p$ on window $w$
under metric $m$, with source and region fixed. Let
$\mathcal W_{p,m}$ be the windows where that value is defined. The
historical aggregate is
\begin{equation}
 \bar y_{p,m}=\frac{1}{|\mathcal W_{p,m}|}
       \sum_{w\in\mathcal W_{p,m}} y_{w,p,m}.
\end{equation}
If this set is empty, the policy has no aggregate for that condition.
Lower-is-better values are negated for cross-metric ranking comparison.
Missing values never become a zero, a worst rank, or a tied rank. A
consequence is that two policy means can describe different observed
window subsets. The pairwise bootstrap difference instead uses the
intersection of windows with both policy outcomes. Agreement of complete
conditions does not remove this conditioning difference.

For comparable policy-value vectors, let $C$ and $D$ be the numbers of
concordant and discordant policy pairs, and $T_x,T_y$ the pairs tied only
in the first or second vector. Kendall's tie-adjusted coefficient is
\begin{equation}
 \tau_b=\frac{C-D}{\sqrt{(C+D+T_x)(C+D+T_y)}}.
\end{equation}
A zero denominator leaves the coefficient undefined. A positive aggregate
value allows individual discordant pairs. Sorting ties lexicographically
for top-$k$ reporting is an implementation convention separate from this
coefficient; it must not silently change the ties counted by tau-b.

For two conditions $x,y$ and policy pair $a,b$, the recorded differences
are $d_x=\bar y_{a,x}-\bar y_{b,x}$ and $d_y=\bar y_{a,y}-\bar y_{b,y}$.
The implementation requires both nonzero and opposite in sign before
considering a reversal. Its relative margin divides the absolute
difference by the absolute aggregate value of the losing policy in that
condition. A zero losing value makes that margin unestimable; it is not
an automatic practical-significance pass. Equal values are assigned zero
difference and zero margin and do not constitute a sign change. These
edge cases matter particularly for bounded utility metrics with many
exact ties.

The saved procedure evaluates eligibility, sign, margin and then interval
support in that order. Consequently, many stable or microscopic records
have no bootstrap interval or $p$-value at all. The BH check is applied to
tests that reach the statistical-support stage within the recorded family,
not to invented $p$-values for every stable record. For the primary
analysis the family is metric/region; for the cross-metric extension it
is source/region/metric-pair. The reported flags are reproducible results
of that procedure. Their interpretation is conditional on this screening,
bootstrap construction and dependence structure; this paper does not prove
unconditional false-discovery control for the entire research workflow.

The saved two-sided bootstrap tail fraction is computed from the same
2,000 resamples as the interval, and can equal zero at finite resolution.
This is not evidence of zero probability of an incorrect scientific claim.
Cross-source ranking bootstraps draw the two source window collections
separately. Cross-metric bootstraps use a common sampled window sequence
because both metrics describe the same outcomes. Repetition IDs are never
resampled as additional windows. M1 verifies saved intervals' classification
logic and point estimates, but does not rerun these stochastic resampling
calculations or substitute new interval estimates for the frozen ones.

\subsection{Evidence denominators}
\begin{table}[t]\centering
\caption{Units in the historical study. Numerically equal totals need not
refer to the same evidence.}\label{tab:denominators}
\begin{tabular}{p{0.43\linewidth}p{0.48\linewidth}}\toprule
Quantity & Interpretation\\\midrule
120 windows & 40 per source; 200 requests in each window\\
9,360 configurations & 120 windows times six regions times 13 policies\\
18,720 executions & Two deterministic passes per configuration\\
18 primary ranking comparisons & Three source pairs times six regions\\
990 primary reversal records & 18 conditions times 55 policy pairs\\
990 cross-metric conditions & Three sources times six regions times 55 metric pairs\\
54,450 metric disagreement records & 990 conditions times 55 policy pairs\\
500 draws per recovery budget & Subsamples of a fixed reference corpus, not new executions\\
240 legacy physical executions & Separate runner and fidelity limits; not included in the simulator matrix\\\bottomrule
\end{tabular}
\end{table}
The inferential unit remains the window. Comparison records reuse its
outcomes many times; their counts describe coverage rather than sample
independence. A change in the number of policy pairs or metrics changes
these denominators without acquiring any additional workload evidence.
